\documentclass[11pt]{article}

\usepackage[margin=1in]{geometry}
\usepackage[T1]{fontenc}
\usepackage[utf8]{inputenc}
\usepackage{lmodern}
\usepackage{microtype}
\usepackage{amsmath,amssymb,amsthm,mathtools}
\usepackage{booktabs}
\usepackage[numbers,sort&compress]{natbib}
\usepackage{hyperref}

\hypersetup{
  colorlinks=true,
  citecolor=blue,
  linkcolor=blue,
  urlcolor=blue
}

\newcommand{\code}[1]{\texttt{#1}}
\newcommand{\bits}[1]{\operatorname{bits}(#1)}
\newcommand{\cost}[1]{\operatorname{cost}(#1)}
\setlength{\emergencystretch}{2em}

\title{A Verified Binary Bit-Cost Reference Model for\\
Integer Normal-Form Algorithms in Lean 4}
\author{Junye Ji\\University of Washington}
\date{Research version 0.1.0, July 22, 2026}

\begin{document}

\maketitle

\begin{abstract}
Formal correctness of a normal-form algorithm does not by itself justify a
bit-complexity claim.  Such a claim must connect abstract integer operations
to a concrete representation, bound all intermediate coefficient widths, and
account for the work of addition, multiplication, division, and gcd.  We
present an executable Lean 4 reference model for canonical binary
sign-magnitude integers.  It provides costed addition with explicit carry and
borrow propagation, schoolbook multiplication, a shared signed Euclidean
quotient/remainder pass, and a well-founded extended Euclidean algorithm.
Each primitive has kernel-checked value correctness, output-width bounds, and
an explicit upper bound on its modeled binary work.  Extended gcd additionally
has a path-sensitive bound for both B\'ezout coefficients and a composed cost
bound driven by the Euclidean quotient sequence.  The development fixes a
four-level vocabulary separating semantic correctness, ring-operation counts,
coefficient growth, and bit complexity.  A reproducible 13-case native corpus
checks execution, but native timing is kept outside the formal claim.  This
model is an independent research component of \code{lean-normal-forms}; it
does not yet assert a bit-complexity theorem for Hermite or Smith normal form.
\end{abstract}

\section{Motivation and scope}

Hermite and Smith normal forms combine matrix transformations with repeated
integer division and gcd.  A proof that the result is canonical answers a
semantic question.  Counting calls to ring operations answers a different
question, since the cost of an integer multiplication depends on operand
size.  Controlling final output size is also insufficient: intermediate
coefficients can be much larger.  A bit-level analysis must instead define a
concrete representation and charge the algorithms that operate on it.
Matrix-level invariants must then supply operand-width bounds before those
charges can be composed.

The \code{lean-normal-forms} core already proves executable HNF and SNF with
strong transformation certificates.  The present component supplies a lower
arithmetic layer for later complexity work.  Its scope is deliberately
narrow:
\begin{itemize}
  \item one canonical binary representation of signed integers;
  \item executable reference addition and multiplication;
  \item signed Euclidean quotient/remainder and extended gcd;
  \item exact implementation-specific counters and independently proved
        upper budgets;
  \item coefficient-width bounds sufficient to compose the extended-gcd
        recurrence.
\end{itemize}
It does not count matrix-level calls, bound normal-form intermediates, or
claim polynomial-time HNF/SNF.  Those are separate obligations for the
Kannan--Bachem and modular-HNF research lines.  This boundary prevents a
primitive arithmetic result from being presented as an algorithm-wide
complexity theorem.

\section{Four acceptance levels}

The Lean type \code{AcceptanceTier} records four project terms:
\begin{enumerate}
  \item \emph{semantic correctness}: the computed value meets its
        mathematical specification;
  \item \emph{ring-operation count}: the number of abstract algebraic
        operations is bounded;
  \item \emph{coefficient bit length}: every relevant intermediate value has
        a proved width bound;
  \item \emph{bit complexity}: concrete binary primitives have explicit
        operation budgets.
\end{enumerate}
The enumeration is vocabulary, not an automatic implication chain.  A
matrix algorithm can establish the first two levels while lacking the third;
without coefficient control, substituting a primitive multiplication cost
does not yield a closed input-size bound.  Likewise, a native benchmark is
neither a ring-operation theorem nor a bit-complexity theorem.

Research version 0.1.0 reaches the fourth level for its four integer
primitives.  The extended-gcd proof includes the third-level invariant needed
to compose multiplication and addition costs at each recursive step.

\section{Representation and proof carriers}

The model reuses mathlib's inductive \code{ZNum} from the manifest-pinned
source commit \citep{mathlib4320znum}.  Its zero constructor is
unique; positive and negative constructors carry a \code{PosNum}, an
inductive, nonzero binary magnitude.  We expose it as
\code{SignMagnitude} rather than introducing a second encoding.  The
functions
\[
  \code{value} : \code{SignMagnitude} \to \mathbb{Z},\qquad
  \code{ofInt} : \mathbb{Z} \to \code{SignMagnitude}
\]
are proved inverse.  The function \code{bitLength} returns zero at zero and
the constructor depth of the positive magnitude otherwise.  Its semantic
theorem identifies that quantity with Lean's binary size of the integer's
absolute value.

A computation returns \code{WithCost $\alpha$}, containing a value and a
natural-number count.  This structure does not grant the count any semantic
meaning by itself.  For every public primitive, three theorem families close
the interpretation:
\begin{itemize}
  \item a value theorem relating the result to standard \code{Int}
        arithmetic;
  \item a bit-length theorem controlling the output;
  \item a cost theorem proving that the returned counter is below a public
        budget.
\end{itemize}
The counter follows the implementation path, whereas the budget is stated in
input widths and proved separately.

\section{Addition and schoolbook multiplication}

\subsection{Carry and borrow}

The addition counter mirrors the recursive equations for positive binary
successor, predecessor, addition, and subtraction.  One unit is charged for
each constructor case inspected; recursive carry and terminal borrow are
charged rather than hidden in a unit-cost arithmetic operation.  The signed
wrapper dispatches on zero and sign, using positive addition for equal signs
and positive subtraction for opposite signs.

For operands with widths $\ell$ and $r$, the public budget is
\[
 B_{+}(\ell,r)=1+2(\ell+r+1)^2.
\]
The implementation returns standard integer addition, its result has width
at most $\ell+r+1$, and its exact modeled counter is at most $B_{+}$.  The
quadratic expression is conservative; the purpose of this reference layer is
a transparent compositional theorem rather than a tight circuit count.

\subsection{Shift and add}

Positive multiplication scans the right operand from least significant bit
to most significant bit.  A zero bit shifts the recursive product; a one bit
also performs a costed addition.  Sign handling is charged by the outer
constructor case.  For left and right widths $\ell$ and $r$, the budget is
\[
 B_{\times}(\ell,r)
 =1+r\bigl(1+(2\ell+r+2)^2\bigr).
\]
The value theorem identifies the result with integer multiplication, and the
product width is at most $\ell+r$.  This is explicitly a schoolbook
reference model, following the classical multiple-precision shift-and-add
model \citep[Section 4.3.1]{knuth1997seminumerical}.  It does not assume the
asymptotics of Lean's native arbitrary-precision integer implementation.

\section{A shared Euclidean long-division pass}

The structure \code{DivModResult} stores quotient and remainder together.
The positive kernel mirrors binary long division
\citep[Section 1.4]{brentZimmermann2010}: it scans numerator bits,
updates a partial remainder, conditionally subtracts the divisor, and emits
one quotient bit.  The signed wrapper reproduces Lean 4.32.0's Euclidean \code{/}
and \code{\%} conventions, including negative numerators, negative divisors,
and division by zero \citep{lean4320intDivMod}.

Producing both fields from one pass is important.  An extended-gcd step needs
both $q$ and $r$ from
\[
  a=qb+r.
\]
Charging two independent divisions would describe a different algorithm.
For widths $n=\bits{a}$ and $d=\bits{b}$, the public budget is
\[
 B_{\mathrm{div}}(n,d)
 =8+n(3+2d)+3n+3d.
\]
The quotient has width at most $n+1$.  When $b\ne 0$, the remainder has width
at most $d$, and its unsigned magnitude is strictly smaller than that of
$b$.  The last theorem is both an arithmetic fact and the termination
argument for the costed extended-gcd loop.

The negative-numerator branch merits explicit treatment.  It divides the
predecessor of the unsigned magnitude, increments the quotient magnitude, and
complements the remainder against the divisor.  Its proof establishes exact
agreement with Lean's signed Euclidean convention rather than silently
switching to truncating division.

\section{Extended gcd and compositional bounds}

The function \code{xgcdWithCost} follows the extended Euclidean recurrence
\citep[Section 1.6]{brentZimmermann2010}.
At a nonzero second argument it computes the shared division, recursively
solves $(b,r)$, and transforms recursive coefficients $(s,t)$ into
\[
  (t,\;s-qt).
\]
The update itself uses \code{mulWithCost} and \code{addWithCost}; their exact
counters are accumulated with the division and recursive costs.  At zero it
normalizes the first argument to a nonnegative gcd and supplies the matching
unit coefficient.

The proposition \code{XGCDSpecification} packages the two semantic claims:
\[
 g=\gcd(a,b),\qquad xa+yb=g.
\]
The theorem \code{xgcdWithCost\_spec} proves both, and a separate theorem
states $g\ge 0$.

Composing primitive costs requires a bound on $s$ and $t$ before the update is
charged.  The recursive function \code{xgcdCoefficientBitBound} follows only
the Euclidean quotient path.  At each return it adds the quotient width and a
conservative multiple of the recursive bound.  The proof shows that both
computed coefficients fit this path-sensitive budget.

The function \code{xgcdBitOperationBound} traverses the same path and sums:
\begin{enumerate}
  \item the current shared-division budget;
  \item the recursive extended-gcd budget;
  \item a multiplication budget instantiated with the quotient width and
        recursive coefficient bound;
  \item an addition budget instantiated with the recursive and product width
        bounds.
\end{enumerate}
It never reads the computed B\'ezout coefficients to choose its budget.  The
exact accumulated counter is bounded by the theorem
\begin{center}
\code{xgcdWithCost\_cost\_le}.
\end{center}
This distinction rules out an after-the-fact ``bound'' that merely repackages
observed intermediate sizes.

\section{Lean organization and trust boundary}

The code is built with Lean 4.32.0 and its matching mathlib snapshot
\citep{deMouraUllrich2021Lean4,lean4320source,mathlib4320source}.  The
research facade is named
\begin{center}
\code{NormalForms.Research.BitCost}.
\end{center}
The frozen core facade does not import it.  A negative-import harness checks
the exact unknown identifier diagnostic, so a syntax failure or unrelated
import failure cannot masquerade as isolation.  A compilation-only surface
freezes 49 declarations.

All executable definitions are constructive and contain no noncomputable
branch.  The 21-root proof audit observes exactly
\code{propext}, \code{Quot.sound}, and \code{Classical.choice}.  The choice
dependency is inherited transitively through semantic lemmas about mathlib's
\code{ZNum}, \code{PosNum}, and \code{Nat.size}; it is not used by an
executable branch.  Project-defined axioms, \code{sorryAx}, and
native/compiler reduction axioms are rejected.

This boundary is intentionally different from the strict no-choice audit for
the core executable normal-form roots.  The current research model proves its
semantics by reusing the standard representation's theorem layer.  Its
complexity claims remain ordinary kernel-checked propositions; no native
decision procedure constructs them.

\section{Executable evaluation}

The native target exercises 13 fixed cases: three additions, three
multiplications, three signed division cases, and four extended-gcd cases.
The corpus covers zero, large carry propagation, cancellation, negative
operands, a zero divisor, a zero gcd input, and signed B\'ezout updates.  Every
row checks both the standard mathematical value and the modeled
cost-versus-budget inequality.

\begin{table}[ht]
\centering
\caption{Selected deterministic results from the 0.1.0 corpus.}
\begin{tabular}{@{}llrrr@{}}
\toprule
Operation & Inputs & Result & Modeled cost & Proved budget \\
\midrule
add    & $-37,19$  & $-18$     & 8   & 289 \\
mul    & $-13,11$  & $-143$    & 18  & 789 \\
divmod & $-37,5$   & $(-8,3)$  & 48  & 89 \\
xgcd   & $240,46$  & $g=2$     & 188 & 298384 \\
\bottomrule
\end{tabular}
\end{table}

The benchmark runner performs one warmup and seven measurements and records
raw wall time, maximum RSS, median, inclusive-quartile IQR, the complete case
output, source revision, executable hash, and hardware fingerprint.  On the
captured AMD Ryzen 9 7940H host, the committed run has a 23.3 ms median and a
69,620 KiB median maximum RSS.  These figures characterize a tiny process
invocation and are not used as CI thresholds or complexity evidence.  The
modeled counters and proved budgets, not native time, are the formal result.

The release gate rebuilds the facade and executable, runs the complete
project test suite, checks all 49 frozen names, enforces the 21-root axiom
set, verifies all 13 native cases, reruns the one-plus-seven protocol, compares
deterministic case data with the baseline, and rebuilds this paper.  A pinned
Debian/Lean container profile provides an independent reproduction path.

\section{Relation to normal-form complexity}

Efficient normal-form algorithms require more than primitive arithmetic.
For example, HNF analyses must bound pivot searches, row operations, and
intermediate determinants; SNF introduces additional gcd and divisibility
steps.  Kannan--Bachem bound both the operation count and coefficient lengths
for HNF and SNF \citep{kannanBachem1979}; Domich--Kannan--Trotter use
modulo-determinant arithmetic for polynomial-time HNF
\citep{domichKannanTrotter1987}.  The present model supplies a concrete cost for
the integer leaves of such an analysis, but does not choose a matrix
algorithm or hide its proof obligations.

The planned Kannan--Bachem line therefore has separate acceptance gates:
semantic HNF, HNF ring-operation count, HNF coefficient growth, HNF bit
complexity, followed by the corresponding SNF stages.  At the final stage, a
matrix-level theorem can invoke the primitive value and cost results here.
The modular-HNF and LLL lines can use the same facade while retaining their
own preconditions, intermediates, and artifacts.

\section{Conclusion}

Research version 0.1.0 turns four common ``integer operations'' into explicit
binary programs with verified semantics and costs.  Canonical
sign-magnitude words make operand width concrete; mirrored counters expose
carry, borrow, shift-and-add, and long-division work; the extended-gcd proof
composes those costs only after establishing a recursive coefficient-width
invariant.  The four-tier vocabulary and independent facade keep the claim at
its actual scope.  This is a reusable arithmetic foundation for later
normal-form complexity proofs, not a substitute for them.

\bibliographystyle{plainnat}
\bibliography{../references}

\end{document}
